\section{The RAG lift hierarchy as a grounding architecture diagnostic}

The 8-scenario results reveal a diagnostic structure:

\begin{center}
\small
\begin{tabular}{@{}ll@{}}
\toprule
Context type & $\Runcited$ \\
\midrule
Step-log (complete reasoning trace) & 0\% \\
Structured domain data (hotspots) & 75\% \\
Maps-only baseline & 60--100\% \\
Narrow structured data (gates) & 100\% \\
\bottomrule
\end{tabular}
\end{center}

This is not a function of data volume but of structural alignment between context type and the science
claims the model otherwise generates.

\textbf{Full grounding (step-log):} The planner step-log records skills dispatched, panels selected,
and artifacts referenced; the query is answerable from those references without open-domain science
claims. $\Runcited$ collapses to 0\% because $|S_\mathrm{sci}| = 0$.

\textbf{Partial grounding (hotspot data):} Structured probability data substitutes for some generic
science claims; the remaining 75\% uncited rate is ecology claims the hotspot data does not cover.

\textbf{No grounding (gate data):} Gate data covers statistical test outcomes only; the model still
generates broader ecology claims gate data does not address.

Context type is a grounding affordance that shapes which claims become evidence-bound~\cite{davis2023affordances};
complete reasoning traces eliminate uncited claims; narrow structured context does not. The diagnostic
ports to any domain where the reasoning trace can be injected as context; eRAG retrieval
evaluation~\cite{arxiv240413781} would show the same pattern --- only step-log documents are fully utilized.

The Citation-Enhanced Generation framework~\cite{arxiv240216063} demonstrates that post-hoc citation
binding reduces hallucination in chatbots. The hierarchy finding complements this: pre-hoc trace
injection eliminates the class of queries that citation binding must address.
